\documentclass[11pt]{article}
\usepackage[margin=1in]{geometry}
\usepackage[utf8]{inputenc}
\usepackage{hyperref}
\usepackage{amsmath,amssymb}
\usepackage{fancyhdr}
\usepackage{enumitem}
\pagestyle{fancy}
\fancyhf{}
\fancyfoot[C]{\thepage}
\title{Configuration spaces and braid groups}
\author{Fred Cohen \and Jonathan Pakianathan}
\date{}
\begin{document}
\maketitle

\noindent\textbf{Source:} \href{https://arxiv.org/abs/2606.17193}{arXiv:2606.17193} \quad \textbf{Field:} math.AT\\

\section*{TL;DR}
The paper presents a survey-style treatment of configuration spaces, braid groups, and mapping class groups on surfaces, unifying K(\ensuremath{\pi},1) characterization, polyfree group theory, and cohomological computations. Specialists converge on a picture of a mathematically sound but underspecified manuscript: novelty is rated significant (0.75, confidence 0.90), individual technical arguments are mostly correct (overall\_correctness: mostly\_sound), and citation coverage is broad (confidence 0.95). However, two major and one critical specialist finding trigger the recommendation gate for the math.* field. The technical\_correctness specialist flagged C8 (major) --- no formal proof artifact for any core induction argument --- and C10 (major) --- the headline genus-2 mapping class group/Sp(4,Z) cohomology claim is asserted in the Introduction but unsubstantiated in the reviewed text. The reproducibility specialist independently flagged a critical concern that the K(\ensuremath{\pi},1), Borel construction, braid group, and mapping class group claims are supported only by informal text with no machine-checkable artifact. These converging findings, combined with several load-bearing arguments that are partially supported (C3, C6, C9) and multiple loop-space and Lie algebra sections that defer to submitted or in-preparation references, require substantive manuscript revision before the paper can be accepted. No specialist disagreements were observed; the technical\_correctness low confidence (0.60) reflects reviewer uncertainty about the uninspected sections (e.g., Section 15), not conflicting assessments.

\textbf{Recommendation:} Major revision \quad \textbf{Confidence:} 72\%\\

\section*{Strengths}
\begin{itemize}[leftmargin=*]
  \item The K(\ensuremath{\pi},1) induction (Theorem 7) using the Fadell-Neuwirth fibration (Theorem 5) is well-executed, with the long exact sequence argument closing correctly for all closed surfaces of genus \ensuremath{\ge} 1 and for the exceptional sphere and projective plane cases handled via Borel constructions.
  \item The polyfree group theory (Propositions 43, 45, 48) is supported by sound LHS spectral-sequence arguments and Nielsen-Schreier induction, yielding the complete package cd(G)=n, type FL, finite presentation, trivial center, and the Euler characteristic product formula.
  \item Novelty is significant: the extension of Fadell-Neuwirth K(\ensuremath{\pi},1) results to the 2-sphere and real projective plane via Borel constructions for SO(3) and S\textsuperscript{3} actions, and the orbit configuration space approach from Xicotencatl, represent nontrivial advances over the prior literature.
  \item Citation coverage is comprehensive and well-integrated, connecting foundational texts (Fadell-Neuwirth, Kohno, Falk-Randell, Cohen-Lada-May) with specialized contemporaneous work (Xicotencatl, D.C. Cohen-Suciu, D.C. Cohen) at high bibliographic confidence.
  \item The paper serves a genuine expository need, assembling scattered results on polyfree groups and their cohomological invariants into a single framework with explicit Euler characteristic formulas.
\end{itemize}

\section*{Weaknesses}
\begin{itemize}[leftmargin=*]
  \item The headline claim that torsion in H*(MCG\_2; Z) and the integral cohomology of Sp(4, Z) 'follows at once' from a configuration-like model (Introduction and Section 15) is asserted without explicit chain-level computation, torsion subgroup generators, or computer-algebra verification, leaving the central advertised application unsubstantiated (C10, major).
  \item No formal proof artifact is provided for any of the load-bearing induction arguments --- the K(\ensuremath{\pi},1) consequence of Fadell-Neuwirth, the polyfree exponent recursion, or the LHS spectral sequence cohomological-dimension lower bound --- in a paper whose field (math.*) places it in the code-amenable set requiring such verification (C8, major; reproducibility concern 2, critical).
  \item Several bibliography entries are critically incomplete: BM (Bergau-Mennicke) and BH (Birman-Hilden) lack title, venue, and year; Smale's entry provides only the venue abbreviation 'P.A.M.S.' with no title or year; and the Xia entry consists solely of the author's name --- all four are cited for substantive results.
  \item The polyfree exponent formula in Theorem 40 defers index verification to the reader as 'an easy exercise,' and the g = 1/2 convention for M = R\textsuperscript{2} is asserted without derivation, creating an off-by-one ambiguity with the series length stated in Remark 41 (C3, C9, minor).
  \item The Lie algebra and loop-space results in the sections relying on CG, CS, and CX are supported only by terse proof sketches and pointers to submitted or in-preparation work, with no standalone derivation sufficient for independent verification (reproducibility concern 4, major).
  \item Arnold (1969) --- the foundational paper establishing the cohomology ring of the pure braid group --- and Birman (1974) --- the standard monograph on braids, links, and mapping class groups --- are absent from the bibliography despite being essential background for the paper's core topics.
\end{itemize}

\section*{Revision Targets}
\begin{itemize}[leftmargin=*]
  \item [open] paper\_tex `corrections/2606.17193/paper.tex` at `Introduction (and Section 15 per the TOC)`: Either provide explicit chain-level computations of H\^{}*(MCG\_2; Z) in Section 15 with explicit torsion subgroup generators, or ship a computer-algebra script (experiments/section15/mcg2\_cohomology.py using Sage's group cohomology package, or a GAP script with HAP) that reproduces the claimed torsion structure and the consequent Sp(4, Z) cohomology computation. Check: Re-review should confirm `Introduction (and Section 15 per the TOC)` is corrected or justified.
  \item [open] paper\_tex `corrections/2606.17193/paper.tex` at `Whole paper`: Ship a formalization repository with at least the core induction lemmas: (i) Fadell-Neuwirth fibration consequence for K(\ensuremath{\pi},1) (src/proofs/FNKpi1.lean), (ii) polyfree exponent recursion (src/proofs/PolyfreeExponents.lean), (iii) the LHS-spectral-sequence cohomological-dimension lower bound (src/proofs/PolyfreeCD.lean). Mathlib already contains the relevant spectral-sequence and group-cohomology infrastructure. Check: Re-review should confirm `Whole paper` is corrected or justified.
  \item [open] bibliography at `J.S. Birman, Braids, Links and Mapping Class Groups, Annals of Mathematics Studies 82, Princeton University Press, 1974`: Add a bibliography entry for `J.S. Birman, Braids, Links and Mapping Class Groups, Annals of Mathematics Studies 82, Princeton University Press, 1974` and cite it where the affected method or claim is introduced, or explicitly justify its omission. Check: Re-review should confirm the bibliography and citation context address this reference.
  \item [open] paper\_tex `corrections/2606.17193/paper.tex` at `Section 5 (Polyfree groups), Theorem 40`: Provide a fully worked verification of the exponent formula in at least one nontrivial case (e.g., g=1, m=0, k=3) and explicitly handle M = R\^{}2 as a separate base case rather than via the g = 1/2 convention. Alternatively, ship a Lean/Coq formalization of the recursion at src/proofs/PolyfreeExponents.lean to mechanically verify the indexing. Check: Re-review should confirm `Section 5 (Polyfree groups), Theorem 40` is corrected or justified.
  \item [open] paper\_tex `corrections/2606.17193/paper.tex` at `reproducibility appendix`: Add a reproducibility note that resolves this concern: Several later Lie algebra and loop-space results rely on terse proof sketches and references to submitted or in-preparation work, without standalone formal verification; an artifact such as proofs/lie\_algebra\_relations.lean would be needed to reproduce these derivations. Check: Re-review should confirm the reproducibility concern is addressed with a concrete artifact or manuscript update.
  \item [open] bibliography at `J.S. Birman, Braids, Links and Mapping Class Groups, Annals of Mathematics Studies 82, Princeton University Press, 1974`: Add a bibliography entry for `J.S. Birman, Braids, Links and Mapping Class Groups, Annals of Mathematics Studies 82, Princeton University Press, 1974` and cite it where the affected method or claim is introduced, or explicitly justify its omission. Check: Re-review should confirm the bibliography and citation context address this reference.
\end{itemize}

\section*{Open Questions}
\begin{itemize}[leftmargin=*]
  \item Can explicit chain-level computations of H*(MCG\_2; Z) be supplied in Section 15, with generators of the torsion subgroup identified, or alternatively a computer-algebra script (Sage/GAP with HAP) reproducing the claimed torsion structure and its consequence for Sp(4, Z)?
  \item Is a formalization of the core induction arguments --- Fadell-Neuwirth K(\ensuremath{\pi},1) consequence, polyfree exponent recursion, and the LHS spectral sequence cohomological-dimension lower bound --- planned or feasible using Lean 4/Mathlib, given that the relevant spectral sequence and group cohomology infrastructure is already available there?
  \item Can the g = 1/2 convention for M = R\textsuperscript{2} in Theorem 40 be either proved as a separate lemma or replaced by an explicit case split, and can the exponent formula be verified in at least one worked example (e.g., g=1, m=0, k=3) to resolve the off-by-one ambiguity with Remark 41?
  \item Will the submitted or in-preparation manuscripts CG (Cohen-Gitler), CS (Cohen-Sato), and CX (Cohen-Xicotencatl) become available to readers who wish to verify the loop-space and orbit configuration space results, or can self-contained proofs replace those deferred arguments in the paper itself?
  \item Can complete bibliographic information be supplied for BM (Bergau-Mennicke), BH (Birman-Hilden), Smale [P.A.M.S.], and Xia, and can Arnold (1969) and Birman (1974) be added to the bibliography?
\end{itemize}

\section*{Per-Agent Reviews}
\subsection*{citation (\texttt{gemini-3-flash-preview})}
\begin{quote}\small\ttfamily\raggedright
\{\\
  "confidence": 0.95,\\
  "entries": [\\
    \{\\
      "citation": \{\\
        "arxiv\_id": null,\\
        "authors": [\\
          "Bredon"\\
        ],\\
        "doi": null,\\
        "key": "Bred",\\
        "raw": "Bred: \{Bredon,\} Topology and Geometry, Springer Verlag GTM 139, New York-Heidelberg-Berlin, 1993.",\\
        "title": "Topology and Geometry",\\
        "url": null,\\
        "venue": "Springer Verlag GTM 139",\\
        "year": 1993\\
      \},\\
      "exists": null,\\
      "explanation": "Standard reference for foundational topology, specifically cited for proofs regarding principal bundles, Thom constructions, and the Künneth theorem.",\\
      "notes": null,\\
      "relevance": "high",\\
      "resolved\_doi": null,\\
      "resolved\_url": null\\
    \},\\
    \{\\
      "citation": \{\\
        "arxiv\_id": null,\\
        "authors": [\\
          "K. S. Brown"\\
        ],\\
        "doi": null,\\
        "key": "Brown",\\
        "raw": "Brown: \{K. S. Brown,\} Cohomology of Groups, Springer Verlag GTM 87, New York-Heidelberg-Berlin, 1994.",\\
        "title": "Cohomology of Groups",\\
        "url": null,\\
        "venue": "Springer Verlag GTM 87",\\
        "year": 1994\\
      \},\\
      "exists": null,\\
      "explanation": "Essential reference for the cohomology of groups, used to bridge topological and algebraic definitions of cohomological dimension.",\\
      "notes": null,\\
      "relevance": "high",\\
      "resolved\_doi": null,\\
      "resolved\_url": null\\
    \},\\
    \{\\
      "citation": \{\\
        "arxiv\_id": null,\\
        "authors": [\\
          "Bergau",\\
          "Mennicke"\\
        ],\\
        "doi": null,\\
        "key": "BM",\\
        "raw": "BM: \{Bergau, and Mennicke\}",\\
        "title": null,\\
        "url": null,\\
        "venue": null,\\
        "year": null\\
      \},\\
      "exists": null,\\
      "explanation": "Presumably supports the discussion on braid groups or mapping class groups, although not explicitly detailed in the provided contexts.",\\
      "notes": "Bibliography entry is incomplete, lacking title, venue, and year.",\\
      "relevance": "medium",\\
      "resolved\_doi": null,\\
      "resolved\_url": null\\
    \},\\
    \{\\
      "citation": \{\\
        "arxiv\_id": null,\\
        "authors": [\\
          "Birman",\\
          "H. Hilden"\\
        ],\\
        "doi": null,\\
        "key": "BH",\\
        "raw": "BH: \{Birman, and H. Hilden,\}",\\
        "title": null,\\
        "url": null,\\
        "venue": null,\\
        "year": null\\
      \},\\
      "exists": null,\\
      "explanation": "Cited for classical methods in the cohomological analysis of mapping class groups.",\\
      "notes": "Bibliography entry is incomplete, lacking title, venue, and year.",\\
      "relevance": "high",\\
      "resolved\_doi": null,\\
      "resolved\_url": null\\
    \},\\
    \{\\
      "citation": \{\\
        "arxiv\_id": null,\\
        "authors": [\\
          "F. R. Cohen"\\
        ],\\
        "doi": null,\\
        "key": "C1",\\
        "raw": "C1: \{F. R. Cohen,\} On the hyperelliptic mapping class groups, \$SO(3)\$, and Spin\$\^{}c(3)\$, American J. Math., \{\textbackslash{}\textbackslash{}bf 115\} (1993), 389--434.",\\
        "title": "On the hyperelliptic mapping class groups, \$SO(3)\$, and Spin\$\^{}c(3)\$ outdoor",\\
        "url": null,\\
        "venue": "American J. Math.",\\
        "year": 1993\\
      \},\\
      "exists": null,\\
      "explanation": "Author's prior work on mapping class groups and SO(3) actions, highly relevant to the paper's focus on the 2-sphere.",\\
      "notes": null,\\
      "relevance": "high",\\
      "resolved\_doi": null,\\
      "resolved\_url": null\\
    \},\\
    \{\\
      "citation": \{\\
        "arxiv\_id": null,\\
        "authors": [\\
          "F. R. Cohen"\\
        ],\\
        "doi": null,\\
        "key": "C2",\\
        "raw": "C2: \{F.R. Cohen,\} On configuration spaces, their homology, and Lie Algebras, Journal of Pure and Applied Algebra \{\textbackslash{}\textbackslash{}bf 100\} (1995), 19-42.",\\
        "title": "On configuration spaces, their homology, and Lie Algebras",\\
        "url": null,\\
        "venue": "Journal of Pure and Applied Algebra",\\
        "year": 1995\\
      \},\\
      "exists": null,\\
      "explanation": "Primary reference by the author for the homological analysis of configuration spaces and associated Lie algebras.",\\
      "notes": null,\\
      "relevance": "high",\\
      "resolved\_doi": null,\\
      "resolved\_url": null\\
    \},\\
    \{\\
      "citation": \{\\
        "arxiv\_id": null,\\
        "authors": [\\
          "F. R. Cohen",\\
          "S. Gitler"\\
        ],\\
        "doi": null,\\
        "key": "CG",\\
        "raw": "CG: \{F.R. Cohen, S. Gitler,\} (i) \{On loop spaces of configuration spaces\}, submitted to the Memoirs of the AMS, and (ii) \{Loop spaces of configuration spaces, braid-like groups, and knots\}, to appear in the Proceedings of the 1998 Barcelona Conference on Algebraic Topology.",\\
        "title": "On loop spaces of configuration spaces / Loop spaces of configuration spaces, braid-like groups, and knots",\\
        "url": null,\\
        "venue": "Submitted / To appear",\\
        "year": 1998\\
      \},\\
      "exists": null,\\
      "explanation": "Key contemporary reference for the main theorem regarding loop spaces of configuration spaces.",\\
      "notes": "Contains two distinct papers within one entry.",\\
      "relevance": "high",\\
      "resolved\_doi": null,\\
      "resolved\_url": null\\
    \},\\
    \{\\
      "citation": \{\\
        "arxiv\_id": null,\\
        "authors": [\\
          "F. R. Cohen",\\
          "T. J. Lada",\\
          "J. P. May"\\
        ],\\
        "doi": null,\\
        "key": "CLM",\\
        "raw": "CLM: \{F.R. Cohen, T.J. Lada and J.P. May,\} The homology of iterated loop spaces, Lecture Notes in Math., vol. \{\textbackslash{}\textbackslash{}bf 533\}, Springer-Verlag (1976).",\\
        "title": "The homology of iterated loop spaces",\\
        "url": null,\\
        "venue": "Lecture Notes in Math., vol. 533",\\
        "year": 1976\\
      \},\\
      "exists": null,\\
      "explanation": "Foundational work on the homology of loop spaces, providing the theoretical basis for the paper's advanced homological analysis.",\\
      "notes": null,\\
      "relevance": "high",\\
      "resolved\_doi": null,\\
      "resolved\_url": null\\
    \},\\
    \{\\
      "citation": \{\\
        "arxiv\_id": null,\\
        "authors": [\\
          "F. R. Cohen",\\
          "T. Sato"\\
        ],\\
        "doi": null,\\
        "key": "CS",\\
        "raw": "CS: \{F.R. Cohen, T. Sato,\} in preparation.",\\
        "title": null,\\
        "url": null,\\
        "venue": "In preparation",\\
        "year": null\\
      \},\\
      "exists": null,\\
      "explanation": "Cited for definitions of fundamental cycles in the homology of configuration spaces.",\\
      "notes": "Work in progress.",\\
      "relevance": "high",\\
      "resolved\_doi": null,\\
      "resolved\_url": null\\
    \},\\
    \{\\
      "citation": \{\\
        "arxiv\_id": null,\\
        "authors": [\\
          "F. R. Cohen",\\
          "M. Xicotencatl"\\
        ],\\
        "doi": null,\\
        "key": "CX",\\
        "raw": "CX: \{F.R. Cohen, M. Xicot\textbackslash{}\textbackslash{}'encatl,\} submitted.",\\
        "title": null,\\
        "url": null,\\
        "venue": "Submitted",\\
        "year": null\\
      \},\\
      "exists": null,\\
      "explanation": "Cited for results on orbit configuration spaces for specific group actions.",\\
      "notes": "Submitted work.",\\
      "relevance": "high",\\
      "resolved\_doi": null,\\
      "resolved\_url": null\\
    \},\\
    \{\\
      "citation": \{\\
        "arxiv\_id": null,\\
        "authors": [\\
          "D. C. Cohen"\\
        ],\\
        "doi": null,\\
        "key": "DC",\\
        "raw": "DC: \{D. C. Cohen,\} Monodromy of fibre-type arrangements and orbit configuration spaces, to appear in Forum Mathematicum.",\\
        "title": "Monodromy of fibre-type arrangements and orbit configuration spaces",\\
        "url": null,\\
        "venue": "Forum Mathematicum",\\
        "year": null\\
      \},\\
      "exists": null,\\
      "explanation": "Directly supports the discussion on orbit configuration spaces and fiber-type arrangements.",\\
      "notes": null,\\
      "relevance": "high",\\
      "resolved\_doi": null,\\
      "resolved\_url": null\\
    \},\\
    \{\\
      "citation": \{\\
        "arxiv\_id": null,\\
        "authors": [\\
          "D. C. Cohen",\\
          "Alexander I. Suciu"\\
        ],\\
        "doi": null,\\
        "key": "Dan",\\
        "raw": "Dan: \{D. C. Cohen, Alexander I. Suciu,\} Homology of iterated semidirect products of free groups, J. Pure Appl. Algebra \$\textbackslash{}\textbackslash{}bf\{126\}\$ (1998), pg 87-120.",\\
        "title": "Homology of iterated semidirect products of free groups",\\
        "url": null,\\
        "venue": "J. Pure Appl. Algebra",\\
        "year": 1998\\
      \},\\
      "exists": null,\\
      "explanation": "Cited for constructions of resolutions for polyfree groups and specific exponents of polyfree series.",\\
      "notes": null,\\
      "relevance": "high",\\
      "resolved\_doi": null,\\
      "resolved\_url": null\\
    \},\\
    \{\\
      "citation": \{\\
        "arxiv\_id": null,\\
        "authors": [\\
          "C. Earle",\\
          "J. Eells"\\
        ],\\
        "doi": null,\\
        "key": "EE",\\
        "raw": "EE: \{C. Earle, and J. Eells,\} A fibre bundle description of Teichm\textbackslash{}\textbackslash{}\textbackslash{}"uller theory, J. Differential Geometry \{\textbackslash{}\textbackslash{}bf 3\}(1969), 19-43.",\\
        "title": "A fibre bundle description of Teichmüller theory",\\
        "url": null,\\
        "venue": "J. Differential Geometry",\\
        "year": 1969\\
      \},\\
      "exists": null,\\
      "explanation": "Used for the homotopy type of diffeomorphism groups and Teichmüller theory in the context of mapping class groups.",\\
      "notes": null,\\
      "relevance": "high",\\
      "resolved\_doi": null,\\
      "resolved\_url": null\\
    \},\\
    \{\\
      "citation": \{\\
        "arxiv\_id": null,\\
        "authors": [\\
          "E. Fadell",\\
          "S. Husseini"\\
        ],\\
        "doi": null,\\
        "key": "FH1",\\
        "raw": "FH1: \{E. Fadell and S. Husseini,\} The space of loops on configuration spaces and the Majer-Terracini index, Topol. Methods in Nonlinear Anal., Journal of the Julius Schauder Center, \{\textbackslash{}\textbackslash{}bf 11\} (1998), 249--271.",\\
        "title": "The space of loops on configuration spaces and the Majer-Terracini index",\\
        "url": null,\\
        "venue": "Topol. Methods in Nonlinear Anal.",\\
        "year": 1998\\
      \},\\
      "exists": null,\\
      "explanation": "Primary source for theorems relating loop spaces of configuration spaces to Lie algebras.",\\
      "notes": null,\\
      "relevance": "high",\\
      "resolved\_doi": null,\\
      "resolved\_url": null\\
    \},\\
    \{\\
      "citation": \{\\
        "arxiv\_id": null,\\
        "authors": [\\
          "E. Fadell",\\
          "S. Husseini"\\
        ],\\
        "doi": null,\\
        "key": "FH2",\\
        "raw": "FH2: \{E. Fadell and S. Husseini,\} Geometry and Topology of Configuration Spaces, in preparation.",\\
        "title": "Geometry and Topology of Configuration Spaces",\\
        "url": null,\\
        "venue": "In preparation",\\
        "year": null\\
      \},\\
      "exists": null,\\
      "explanation": "Supporting reference for the broader theory of configuration spaces developed by Fadell and Husseini.",\\
      "notes": "Work in progress.",\\
      "relevance": "high",\\
      "resolved\_doi": null,\\
      "resolved\_url": null\\
    \},\\
    \{\\
      "citation": \{\\
        "arxiv\_id": null,\\
        "authors": [\\
          "E. Fadell",\\
          "L. Neuwirth"\\
        ],\\
        "doi": null,\\
        "key": "FN",\\
        "raw": "FN: \{E. Fadell and L. Neuwirth,\} Configuration spaces, Math. Scand. \{\textbackslash{}\textbackslash{}bf 10\} (1962), 119-126.",\\
        "title": "Configuration spaces",\\
        "url": null,\\
        "venue": "Math. Scand.",\\
        "year": 1962\\
      \},\\
      "exists": null,\\
      "explanation": "The seminal paper defining configuration spaces and their basic properties, essential for any work in the field.",\\
      "notes": null,\\
      "relevance": "high",\\
      "resolved\_doi": null,\\
      "resolved\_url": null\\
    \},\\
    \{\\
      "citation": \{\\
        "arxiv\_id": null,\\
        "authors": [\\
          "M. Falk",\\
          "R. Randell"\\
        ],\\
        "doi": null,\\
        "key": "FR",\\
        "raw": "FR: \{M. Falk, and R. Randell,\} The lower central series of a fiber-type arrangement, Invent. Math. \{\textbackslash{}\textbackslash{}bf 82\} (1985), 77-88.",\\
        "title": "The lower central series of a fiber-type arrangement",\\
        "url": null,\\
        "venue": "Invent. Math.",\\
        "year": 1985\\
      \},\\
      "exists": null,\\
      "explanation": "Foundational for the Kohno-Falk-Randell theory described in the paper regarding Lie algebras of braid groups.",\\
      "notes": null,\\
      "relevance": "high",\\
      "resolved\_doi": null,\\
      "resolved\_url": null\\
    \},\\
    \{\\
      "citation": \{\\
        "arxiv\_id": null,\\
        "authors": [\\
          "V. Guillemin",\\
          "A. Pollack"\\
        ],\\
        "doi": null,\\
        "key": "GP",\\
        "raw": "GP: \{V. Guillemin, and A. Pollack,\} Differential Topology, Prentice-Hall, 1974.",\\
        "title": "Differential Topology",\\
        "url": null,\\
        "venue": "Prentice-Hall",\\
        "year": 1974\\
      \},\\
      "exists": null,\\
      "explanation": "Cited for the standard isotopy lemma used in the section on mapping class groups.",\\
      "notes": null,\\
      "relevance": "high",\\
      "resolved\_doi": null,\\
      "resolved\_url": null\\
    \},\\
    \{\\
      "citation": \{\\
        "arxiv\_id": null,\\
        "authors": [\\
          "R. Hain"\\
        ],\\
        "doi": null,\\
        "key": "H",\\
        "raw": "H: \{R. Hain,\} Infinitesimal presentations of the Torelli group, J. Amer. Math. Soc. \{\textbackslash{}\textbackslash{}bf 10\} (1997),597-691.",\\
        "title": "Infinitesimal presentations of the Torelli group",\\
        "url": null,\\
        "venue": "J. Amer. Math. Soc.",\\
        "year": 1997\\
      \},\\
      "exists": null,\\
      "explanation": "Relevant to the homological analysis of mapping class groups, specifically the Torelli subgroup.",\\
      "notes": null,\\
      "relevance": "high",\\
      "resolved\_doi": null,\\
      "resolved\_url": null\\
    \},\\
    \{\\
      "citation": \{\\
        "arxiv\_id": null,\\
        "authors": [\\
          "T. Kohno"\\
        ],\\
        "doi": null,\\
        "key": "K1",\\
        "raw": "K1: \{T. Kohno,\} Linear represenations of braid groups and classical Yang-Baxter equations, Cont. Math. \{\textbackslash{}\textbackslash{}bf 78\} (1988), 339-363.",\\
        "title": "Linear represenations of braid groups and classical Yang-Baxter equations",\\
        "url": null,\\
        "venue": "Cont. Math.",\\
        "year": 1988\\
      \},\\
      "exists": null,\\
      "explanation": "Part of Kohno's extensive work on braid groups, specifically relating to Lie algebras.",\\
      "notes": null,\\
      "relevance": "high",\\
      "resolved\_doi": null,\\
      "resolved\_url": null\\
    \},\\
    \{\\
      "citation": \{\\
        "arxiv\_id": null,\\
        "authors": [\\
          "T. Kohno"\\
        ],\\
        "doi": null,\\
        "key": "K2",\\
        "raw": "K2: \{T. Kohno,\} Vassiliev invariants and de Rham complex on the space of knots, Cont. Math. \{\textbackslash{}\textbackslash{}bf 179\} (1994), 123-138.",\\
        "title": "Vassiliev invariants and de Rham complex on the space of knots",\\
        "url": null,\\
        "venue": "Cont. Math.",\\
        "year": 1994\\
      \},\\
      "exists": null,\\
      "explanation": "Relates configuration spaces to knot invariants, a context mentioned in the text.",\\
      "notes": null,\\
      "relevance": "high",\\
      "resolved\_doi": null,\\
      "resolved\_url": null\\
    \},\\
    \{\\
      "citation": \{\\
        "arxiv\_id": null,\\
        "authors": [\\
          "T. Kohno"\\
        ],\\
        "doi": null,\\
        "key": "K3",\\
        "raw": "K3: \{T. Kohno,\} Elliptic KZ system, braid groups of the torus and Vassiliev invariants, Topology and its Applications, \{\textbackslash{}\textbackslash{}bf 78\} (1997), 79-94.",\\
        "title": "Elliptic KZ system, braid groups of the torus and Vassiliev invariants",\\
        "url": null,\\
        "venue": "Topology and its Applications",\\
        "year": 1997\\
      \},\\
      "exists": null,\\
      "explanation": "Directly relevant to the discussion of braids on the torus and related invariants.",\\
      "notes": null,\\
      "relevance": "high",\\
      "resolved\_doi": null,\\
      "resolved\_url": null\\
    \},\\
    \{\\
      "citation": \{\\
        "arxiv\_id": null,\\
        "authors": [\\
          "T. Kohno",\\
          "T. Oda"\\
        ],\\
        "doi": null,\\
        "key": "KO",\\
        "raw": "KO: \{T. Kohno, and T. Oda,\} The lower central series of the pure braid group of an algebraic curve, Advanced Studies in Pure Mathematics, Galois Represntations and Arithmetic Algebraic Geometry \{\textbackslash{}\textbackslash{}bf 12\} (1987), 201-219.",\\
        "title": "The lower central series of the pure braid group of an algebraic curve",\\
        "url": null,\\
        "venue": "Advanced Studies in Pure Mathematics",\\
        "year": 1987\\
      \},\\
      "exists": null,\\
      "explanation": "Cited for the analysis of the descending central series of pure braid groups of curves.",\\
      "notes": null,\\
      "relevance": "high",\\
      "resolved\_doi": null,\\
      "resolved\_url": null\\
    \},\\
    \{\\
      "citation": \{\\
        "arxiv\_id": null,\\
        "authors": [\\
          "J. W. Milnor",\\
          "J. C. Moore"\\
        ],\\
        "doi": null,\\
        "key": "MM",\\
        "raw": "MM: \{J. W. Milnor, and J.C. Moore,\} On the structure of Hopf algebras, Ann. of Math. \{\textbackslash{}\textbackslash{}bf 81\} (1965), 211-264.",\\
        "title": "On the structure of Hopf algebras",\\
        "url": null,\\
        "venue": "Ann. of Math.",\\
        "year": 1965\\
      \},\\
      "exists": null,\\
      "explanation": "Foundational reference for Hopf algebras, necessary for the homological treatment of loop spaces.",\\
      "notes": null,\\
      "relevance": "high",\\
      "resolved\_doi": null,\\
      "resolved\_url": null\\
    \},\\
    \{\\
      "citation": \{\\
        "arxiv\_id": null,\\
        "authors": [\\
          "Derek J. S. Robinson"\\
        ],\\
        "doi": null,\\
        "key": "Robin",\\
        "raw": "Robin: \{Derek J.S. Robinson,\} A Course in the Theory of Groups, Springer Verlag GTM 80, New York-Heidelberg-Berlin, 1995.",\\
        "title": "A Course in the Theory of Groups",\\
        "url": null,\\
        "venue": "Springer Verlag GTM 80",\\
        "year": 1995\\
      \},\\
      "exists": null,\\
      "explanation": "Standard group theory text cited for the Nielsen-Schreier theorem in the context of polyfree groups.",\\
      "notes": null,\\
      "relevance": "high",\\
      "resolved\_doi": null,\\
      "resolved\_url": null\\
    \},\\
    \{\\
      "citation": \{\\
        "arxiv\_id": null,\\
        "authors": [\\
          "N. E. Steenrod"\\
        ],\\
        "doi": null,\\
        "key": "S",\\
        "raw": "S: \{N. E. Steenrod,\} The Topology of Fibre Bundles, Princeton University Press.",\\
        "title": "The Topology of Fibre Bundles",\\
        "url": null,\\
        "venue": "Princeton University Press",\\
        "year": null\\
      \},\\
      "exists": null,\\
      "explanation": "Classic text on fiber bundles, cited for specific foundational results.",\\
      "notes": "Missing publication year.",\\
      "relevance": "high",\\
      "resolved\_doi": null,\\
      "resolved\_url": null\\
    \},\\
    \{\\
      "citation": \{\\
        "arxiv\_id": null,\\
        "authors": [\\
          "T. Sato"\\
        ],\\
        "doi": null,\\
        "key": "Sa",\\
        "raw": "Sa: \{T. Sato,\} On free groups and morphisms of coalgebras, student topology seminar, Univ. of Roch., Spring 1998, preprint.",\\
        "title": "On free groups and morphisms of coalgebras",\\
        "url": null,\\
        "venue": "Univ. of Roch., student topology seminar",\\
        "year": 1998\\
      \},\\
      "exists": null,\\
      "explanation": "Seminar notes supporting the homological algebra used in the paper.",\\
      "notes": "Student seminar preprint.",\\
      "relevance": "medium",\\
      "resolved\_doi": null,\\
      "resolved\_url": null\\
    \},\\
    \{\\
      "citation": \{\\
        "arxiv\_id": null,\\
        "authors": [\\
          "S. Smale"\\
        ],\\
        "doi": null,\\
        "key": "Smale",\\
        "raw": "Smale: \{S. Smale,\} P.A.M.S.",\\
        "title": null,\\
        "url": null,\\
        "venue": "P.A.M.S.",\\
        "year": null\\
      \},\\
      "exists": null,\\
      "explanation": "Cited for the homotopy equivalence between SO(3) and Diff+(S2).",\\
      "notes": "Entry is extremely sparse, missing title and year.",\\
      "relevance": "high",\\
      "resolved\_doi": null,\\
      "resolved\_url": null\\
    \},\\
    \{\\
      "citation": \{\\
        "arxiv\_id": null,\\
        "authors": [\\
          "Edwin H. Spanier"\\
        ],\\
        "doi": null,\\
        "key": "Span",\\
        "raw": "Span: \{Edwin H. Spanier,\} Algebraic Topology, Springer Verlag, New York-Heidelberg-Berlin, 1966.",\\
        "title": "Algebraic Topology",\\
        "url": null,\\
        "venue": "Springer Verlag",\\
        "year": 1966\\
      \},\\
      "exists": null,\\
      "explanation": "Standard algebraic topology text cited for basic properties alongside Bredon.",\\
      "notes": null,\\
      "relevance": "high",\\
      "resolved\_doi": null,\\
      "resolved\_url": null\\
    \},\\
    \{\\
      "citation": \{\\
        "arxiv\_id": null,\\
        "authors": [\\
          "Whitehead"\\
        ],\\
        "doi": null,\\
        "key": "White",\\
        "raw": "White: \{Whitehead,\} Elements of Homotopy Theory, Springer Verlag GTM 61, New York-Heidelberg-Berlin, .",\\
        "title": "Elements of Homotopy Theory",\\
        "url": null,\\
        "venue": "Springer Verlag GTM 61",\\
        "year": null\\
      \},\\
      "exists": null,\\
      "explanation": "Comprehensive reference for homotopy theory, likely supporting the general categorical framework.",\\
      "notes": "Entry ends abruptly with a comma.",\\
      "relevance": "high",\\
      "resolved\_doi": null,\\
      "resolved\_url": null\\
    \},\\
    \{\\
      "citation": \{\\
        "arxiv\_id": null,\\
        "authors": [\\
          "S. Willerton"\\
        ],\\
        "doi": null,\\
        "key": "W",\\
        "raw": "W: \{S. Willerton,\} thesis.",\\
        "title": null,\\
        "url": null,\\
        "venue": "Thesis",\\
        "year": null\\
      \},\\
      "exists": null,\\
      "explanation": "Presumably relates to Vassiliev invariants, a topic mentioned but not detailed in the paper.",\\
      "notes": "Very vague citation.",\\
      "relevance": "medium",\\
      "resolved\_doi": null,\\
      "resolved\_url": null\\
    \},\\
    \{\\
      "citation": \{\\
        "arxiv\_id": null,\\
        "authors": [\\
          "M. Xicotencatl"\\
        ],\\
        "doi": null,\\
        "key": "X",\\
        "raw": "X: \{M. Xicot\textbackslash{}\textbackslash{}'encatl,\} Orbit configuration spaces, infinitesimal braid relations, and equivariant function spaces, Ph.D. thesis, University of Rochester, Spring 1997.",\\
        "title": "Orbit configuration spaces, infinitesimal braid relations, and equivariant function spaces",\\
        "url": null,\\
        "venue": "Ph.D. thesis, University of Rochester",\\
        "year": 1997\\
      \},\\
      "exists": null,\\
      "explanation": "The primary source for the theory of orbit configuration spaces, which is a major focus of the latter half of the paper.",\\
      "notes": null,\\
      "relevance": "high",\\
      "resolved\_doi": null,\\
      "resolved\_url": null\\
    \},\\
    \{\\
      "citation": \{\\
        "arxiv\_id": null,\\
        "authors": [\\
          "Yining Xia"\\
        ],\\
        "doi": null,\\
        "key": "Xia",\\
        "raw": "Xia: \{Yining Xia,\}",\\
        "title": null,\\
        "url": null,\\
        "venue": null,\\
        "year": null\\
      \},\\
      "exists": null,\\
      "explanation": "Presumably supports specific homological results, though context is not provided.",\\
      "notes": "Incomplete entry containing only the author's name.",\\
      "relevance": "medium",\\
      "resolved\_doi": null,\\
      "resolved\_url": null\\
    \}\\
  ],\\
  "missing\_references": [\\
    \{\\
      "reason": "This monograph is the standard and comprehensive reference for the general theory of braids and mapping class groups discussed throughout the paper.",\\
      "title": "J.S. Birman, Braids, Links and Mapping Class Groups, Annals of Mathematics Studies 82, Princeton University Press, 1974"\\
    \},\\
    \{\\
      "reason": "This foundational paper established the cohomology ring of the pure braid group, which is the starting point for much of the homological analysis presented.",\\
      "title": "V.I. Arnold, The cohomology ring of the colored braid group, Mathematical Notes of the Academy of Sciences of the USSR 5 (1969), 138-140"\\
    \}\\
  ],\\
  "summary": "The paper exhibits high citation hygiene, effectively integrating foundational topological texts (Bredon, Spanier, Steenrod) with specialized contemporary research (Xicotencatl, Kohno, Fadell). While the bibliography is comprehensive, several entries (e.g., BM, BH, Xia, Smale) are significantly underdeveloped, lacking titles, years, or venues. The addition of Birman's 1974 monograph would further strengthen the background on braid groups."\\
\}\\
\end{quote}

\subsection*{meta\_reviewer (\texttt{claude-sonnet-4-6})}
\begin{quote}\small\ttfamily\raggedright
\{\\
  "confidence": 0.72,\\
  "questions": [\\
    "Can explicit chain-level computations of H*(MCG\_2; Z) be supplied in Section 15, with generators of the torsion subgroup identified, or alternatively a computer-algebra script (Sage/GAP with HAP) reproducing the claimed torsion structure and its consequence for Sp(4, Z)?",\\
    "Is a formalization of the core induction arguments --- Fadell-Neuwirth K(\ensuremath{\pi},1) consequence, polyfree exponent recursion, and the LHS spectral sequence cohomological-dimension lower bound --- planned or feasible using Lean 4/Mathlib, given that the relevant spectral sequence and group cohomology infrastructure is already available there?",\\
    "Can the g = 1/2 convention for M = R\textsuperscript{2} in Theorem 40 be either proved as a separate lemma or replaced by an explicit case split, and can the exponent formula be verified in at least one worked example (e.g., g=1, m=0, k=3) to resolve the off-by-one ambiguity with Remark 41?",\\
    "Will the submitted or in-preparation manuscripts CG (Cohen-Gitler), CS (Cohen-Sato), and CX (Cohen-Xicotencatl) become available to readers who wish to verify the loop-space and orbit configuration space results, or can self-contained proofs replace those deferred arguments in the paper itself?",\\
    "Can complete bibliographic information be supplied for BM (Bergau-Mennicke), BH (Birman-Hilden), Smale [P.A.M.S.], and Xia, and can Arnold (1969) and Birman (1974) be added to the bibliography?"\\
  ],\\
  "recommendation": "major\_revision",\\
  "revision\_targets": [\\
    \{\\
      "evidence": "This headline claim from the Introduction is asserted but not substantiated in the portion of the paper available for review. The promise that 'torsion in the cohomology follows at once' from a configuration-like model is a strong assertion; without the explicit computation visible in the inspected text, and without an executable cohomology computation (e.g., via Sage, GAP, or a formalized Lean calculation against a CW model), the claim is at most a survey-level pointer.",\\
      "id": "weakness-1",\\
      "locator": "Introduction (and Section 15 per the TOC)",\\
      "required\_update": "Either provide explicit chain-level computations of H\^{}*(MCG\_2; Z) in Section 15 with explicit torsion subgroup generators, or ship a computer-algebra script (experiments/section15/mcg2\_cohomology.py using Sage's group cohomology package, or a GAP script with HAP) that reproduces the claimed torsion structure and the consequent Sp(4, Z) cohomology computation.",\\
      "source\_path": "corrections/2606.17193/paper.tex",\\
      "source\_role": "technical\_correctness",\\
      "status": "open",\\
      "target\_kind": "paper\_tex",\\
      "verification\_check": "Re-review should confirm `Introduction (and Section 15 per the TOC)` is corrected or justified.",\\
      "weakness\_index": 0\\
    \},\\
    \{\\
      "evidence": "Per the Proof-as-Code Axiom for math.* fields: although the paper's individual results are classical or carefully argued, no executable formal verification artifact (Coq, Lean, Agda, or Isabelle/HOL) is shipped to mechanically certify the recursion arguments (Theorem 7, Theorem 40, Proposition 43, Lemma 47, Proposition 48). For a 'notes' / survey style paper that reproves classical results, the absence of a formalization repository is evidence that the load-bearing recursions have not been independently mechanically checked.",\\
      "id": "weakness-2",\\
      "locator": "Whole paper",\\
      "required\_update": "Ship a formalization repository with at least the core induction lemmas: (i) Fadell-Neuwirth fibration consequence for K(\ensuremath{\pi},1) (src/proofs/FNKpi1.lean), (ii) polyfree exponent recursion (src/proofs/PolyfreeExponents.lean), (iii) the LHS-spectral-sequence cohomological-dimension lower bound (src/proofs/PolyfreeCD.lean). Mathlib already contains the relevant spectral-sequence and group-cohomology infrastructure.",\\
      "source\_path": "corrections/2606.17193/paper.tex",\\
      "source\_role": "technical\_correctness",\\
      "status": "open",\\
      "target\_kind": "paper\_tex",\\
      "verification\_check": "Re-review should confirm `Whole paper` is corrected or justified.",\\
      "weakness\_index": 1\\
    \},\\
    \{\\
      "evidence": "This monograph is the standard and comprehensive reference for the general theory of braids and mapping class groups discussed throughout the paper.",\\
      "id": "weakness-3",\\
      "locator": "J.S. Birman, Braids, Links and Mapping Class Groups, Annals of Mathematics Studies 82, Princeton University Press, 1974",\\
      "required\_update": "Add a bibliography entry for `J.S. Birman, Braids, Links and Mapping Class Groups, Annals of Mathematics Studies 82, Princeton University Press, 1974` and cite it where the affected method or claim is introduced, or explicitly justify its omission.",\\
      "source\_path": null,\\
      "source\_role": "citation",\\
      "status": "open",\\
      "target\_kind": "bibliography",\\
      "verification\_check": "Re-review should confirm the bibliography and citation context address this reference.",\\
      "weakness\_index": 2\\
    \},\\
    \{\\
      "evidence": "The induction step uses the Fadell-Neuwirth fibration \ensuremath{\pi}\_\{k,k-1\} with fiber F(M \ensuremath{-} Q\_\{m+k-1\}, 1), correctly yielding a free factor of rank 2g + m \ensuremath{-} 2 + k. Combining this with the inductive series gives a polyfree series of length k. The exponent formula 2g + m \ensuremath{-} 1 \ensuremath{-} i + k matches this at the boundary values (i=k gives 2g+m\ensuremath{-}1, the rank of \ensuremath{\pi}\_1(M \ensuremath{-} Q\_m); i=1 gives 2g+m\ensuremath{-}2+k, the rank of the top fiber added in the inductive step), but the proof explicitly defers the verification of the indexing to the reader ('It is an easy exercise which will be left to the reader'). The unconventional choice g = 1/2 for R\^{}2 is a notational hack that compresses two cases but is not rigorously justified --- the underlying claim \ensuremath{\pi}\_1(R\^{}2 \ensuremath{-} Q\_m) is free of rank m needs separate treatment.",\\
      "id": "weakness-4",\\
      "locator": "Section 5 (Polyfree groups), Theorem 40",\\
      "required\_update": "Provide a fully worked verification of the exponent formula in at least one nontrivial case (e.g., g=1, m=0, k=3) and explicitly handle M = R\^{}2 as a separate base case rather than via the g = 1/2 convention. Alternatively, ship a Lean/Coq formalization of the recursion at src/proofs/PolyfreeExponents.lean to mechanically verify the indexing.",\\
      "source\_path": "corrections/2606.17193/paper.tex",\\
      "source\_role": "technical\_correctness",\\
      "status": "open",\\
      "target\_kind": "paper\_tex",\\
      "verification\_check": "Re-review should confirm `Section 5 (Polyfree groups), Theorem 40` is corrected or justified.",\\
      "weakness\_index": 3\\
    \},\\
    \{\\
      "evidence": "Several later Lie algebra and loop-space results rely on terse proof sketches and references to submitted or in-preparation work, without standalone formal verification; an artifact such as proofs/lie\_algebra\_relations.lean would be needed to reproduce these derivations.",\\
      "id": "weakness-5",\\
      "locator": "reproducibility appendix",\\
      "required\_update": "Add a reproducibility note that resolves this concern: Several later Lie algebra and loop-space results rely on terse proof sketches and references to submitted or in-preparation work, without standalone formal verification; an artifact such as proofs/lie\_algebra\_relations.lean would be needed to reproduce these derivations.",\\
      "source\_path": "corrections/2606.17193/paper.tex",\\
      "source\_role": "reproducibility",\\
      "status": "open",\\
      "target\_kind": "paper\_tex",\\
      "verification\_check": "Re-review should confirm the reproducibility concern is addressed with a concrete artifact or manuscript update.",\\
      "weakness\_index": 4\\
    \},\\
    \{\\
      "evidence": "This monograph is the standard and comprehensive reference for the general theory of braids and mapping class groups discussed throughout the paper.",\\
      "id": "weakness-6",\\
      "locator": "J.S. Birman, Braids, Links and Mapping Class Groups, Annals of Mathematics Studies 82, Princeton University Press, 1974",\\
      "required\_update": "Add a bibliography entry for `J.S. Birman, Braids, Links and Mapping Class Groups, Annals of Mathematics Studies 82, Princeton University Press, 1974` and cite it where the affected method or claim is introduced, or explicitly justify its omission.",\\
      "source\_path": null,\\
      "source\_role": "citation",\\
      "status": "open",\\
      "target\_kind": "bibliography",\\
      "verification\_check": "Re-review should confirm the bibliography and citation context address this reference.",\\
      "weakness\_index": 5\\
    \}\\
  ],\\
  "strengths": [\\
    "The K(\ensuremath{\pi},1) induction (Theorem 7) using the Fadell-Neuwirth fibration (Theorem 5) is well-executed, with the long exact sequence argument closing correctly for all closed surfaces of genus \ensuremath{\ge} 1 and for the exceptional sphere and projective plane cases handled via Borel constructions.",\\
    "The polyfree group theory (Propositions 43, 45, 48) is supported by sound LHS spectral-sequence arguments and Nielsen-Schreier induction, yielding the complete package cd(G)=n, type FL, finite presentation, trivial center, and the Euler characteristic product formula.",\\
    "Novelty is significant: the extension of Fadell-Neuwirth K(\ensuremath{\pi},1) results to the 2-sphere and real projective plane via Borel constructions for SO(3) and S\textsuperscript{3} actions, and the orbit configuration space approach from Xicotencatl, represent nontrivial advances over the prior literature.",\\
    "Citation coverage is comprehensive and well-integrated, connecting foundational texts (Fadell-Neuwirth, Kohno, Falk-Randell, Cohen-Lada-May) with specialized contemporaneous work (Xicotencatl, D.C. Cohen-Suciu, D.C. Cohen) at high bibliographic confidence.",\\
    "The paper serves a genuine expository need, assembling scattered results on polyfree groups and their cohomological invariants into a single framework with explicit Euler characteristic formulas."\\
  ],\\
  "summary": "The paper presents a survey-style treatment of configuration spaces, braid groups, and mapping class groups on surfaces, unifying K(\ensuremath{\pi},1) characterization, polyfree group theory, and cohomological computations. Specialists converge on a picture of a mathematically sound but underspecified manuscript: novelty is rated significant (0.75, confidence 0.90), individual technical arguments are mostly correct (overall\_correctness: mostly\_sound), and citation coverage is broad (confidence 0.95). However, two major and one critical specialist finding trigger the recommendation gate for the math.* field. The technical\_correctness specialist flagged C8 (major) --- no formal proof artifact for any core induction argument --- and C10 (major) --- the headline genus-2 mapping class group/Sp(4,Z) cohomology claim is asserted in the Introduction but unsubstantiated in the reviewed text. The reproducibility specialist independently flagged a critical concern that the K(\ensuremath{\pi},1), Borel construction, braid group, and mapping class group claims are supported only by informal text with no machine-checkable artifact. These converging findings, combined with several load-bearing arguments that are partially supported (C3, C6, C9) and multiple loop-space and Lie algebra sections that defer to submitted or in-preparation references, require substantive manuscript revision before the paper can be accepted. No specialist disagreements were observed; the technical\_correctness low confidence (0.60) reflects reviewer uncertainty about the uninspected sections (e.g., Section 15), not conflicting assessments.",\\
  "weaknesses": [\\
    "The headline claim that torsion in H*(MCG\_2; Z) and the integral cohomology of Sp(4, Z) 'follows at once' from a configuration-like model (Introduction and Section 15) is asserted without explicit chain-level computation, torsion subgroup generators, or computer-algebra verification, leaving the central advertised application unsubstantiated (C10, major).",\\
    "No formal proof artifact is provided for any of the load-bearing induction arguments --- the K(\ensuremath{\pi},1) consequence of Fadell-Neuwirth, the polyfree exponent recursion, or the LHS spectral sequence cohomological-dimension lower bound --- in a paper whose field (math.*) places it in the code-amenable set requiring such verification (C8, major; reproducibility concern 2, critical).",\\
    "Several bibliography entries are critically incomplete: BM (Bergau-Mennicke) and BH (Birman-Hilden) lack title, venue, and year; Smale's entry provides only the venue abbreviation 'P.A.M.S.' with no title or year; and the Xia entry consists solely of the author's name --- all four are cited for substantive results.",\\
    "The polyfree exponent formula in Theorem 40 defers index verification to the reader as 'an easy exercise,' and the g = 1/2 convention for M = R\textsuperscript{2} is asserted without derivation, creating an off-by-one ambiguity with the series length stated in Remark 41 (C3, C9, minor).",\\
    "The Lie algebra and loop-space results in the sections relying on CG, CS, and CX are supported only by terse proof sketches and pointers to submitted or in-preparation work, with no standalone derivation sufficient for independent verification (reproducibility concern 4, major).",\\
    "Arnold (1969) --- the foundational paper establishing the cohomology ring of the pure braid group --- and Birman (1974) --- the standard monograph on braids, links, and mapping class groups --- are absent from the bibliography despite being essential background for the paper's core topics."\\
  ]\\
\}\\
\end{quote}

\subsection*{novelty (\texttt{gemini-3-flash-preview})}
\begin{quote}\small\ttfamily\raggedright
\{\\
  "confidence": 0.9,\\
  "missing\_prior\_art": [],\\
  "novelty\_score": 0.75,\\
  "related\_work": [\\
    \{\\
      "citation\_key": "FN",\\
      "delta": "Extends the Fadell-Neuwirth theorem, which establishes that configuration spaces of manifolds are \$K(\textbackslash{}\textbackslash{}pi,1)\$ spaces for most surfaces, to the exceptional cases of the 2-sphere and the real projective plane by utilizing Borel constructions for \$SO(3)\$ and \$S\^{}3\$ actions.",\\
      "relation": "builds\_on",\\
      "title": "Configuration spaces"\\
    \},\\
    \{\\
      "citation\_key": "X",\\
      "delta": "Incorporates the theory of orbit configuration spaces to analyze the configuration spaces of the real projective plane, establishing their \$K(\textbackslash{}\textbackslash{}pi,1)\$ homotopy type through equivariant models.",\\
      "relation": "builds\_on",\\
      "title": "Orbit configuration spaces, infinitesimal braid relations, and equivariant function spaces"\\
    \},\\
    \{\\
      "citation\_key": "FR",\\
      "delta": "Applies the Kohno-Falk-Randell theory of Lie algebras associated with descending central series to the loop spaces of configuration spaces, providing a homological interpretation of braid group structures.",\\
      "relation": "prior\_art",\\
      "title": "The lower central series of a fiber-type arrangement"\\
    \},\\
    \{\\
      "citation\_key": "Dan",\\
      "delta": "Utilizes specific resolutions and Fox free derivatives for polyfree groups to derive representations and homological properties of pure braid groups.",\\
      "relation": "builds\_on",\\
      "title": "Homology of iterated semidirect products of free groups"\\
    \}\\
  ],\\
  "verdict": "significant"\\
\}\\
\end{quote}

\subsection*{reproducibility (\texttt{gpt-5.5})}
\begin{quote}\small\ttfamily\raggedright
\{\\
  "code\_availability": "unspecified",\\
  "code\_url": null,\\
  "concerns": [\\
    \{\\
      "area": "code",\\
      "description": "No code repository, formal proof repository, commit, release, or license is provided for reproducing the paper's mathematical claims.",\\
      "severity": "major"\\
    \},\\
    \{\\
      "area": "other",\\
      "description": "The headline K(pi,1), Borel construction, braid group, and mapping class group claims are supported only by informal text; a machine-checkable artifact such as proofs/configuration\_spaces\_and\_braid\_groups.lean would be needed to close this gap.",\\
      "severity": "critical"\\
    \},\\
    \{\\
      "area": "other",\\
      "description": "The homological and cohomological calculations for genus zero, one, and two surfaces are not accompanied by formalized computations or checkable tables; an artifact such as proofs/cohomology\_calculations.lean or scripts/verify\_cohomology.sage would be needed for independent reproduction.",\\
      "severity": "major"\\
    \},\\
    \{\\
      "area": "other",\\
      "description": "Several later Lie algebra and loop-space results rely on terse proof sketches and references to submitted or in-preparation work, without standalone formal verification; an artifact such as proofs/lie\_algebra\_relations.lean would be needed to reproduce these derivations.",\\
      "severity": "major"\\
    \}\\
  ],\\
  "confidence": 0.86,\\
  "data\_availability": "unspecified",\\
  "data\_url": null,\\
  "environment": \{\\
    "dependencies": [],\\
    "hardware": null,\\
    "software": null\\
  \},\\
  "reproducibility\_score": 0.22\\
\}\\
\end{quote}

\subsection*{summary (\texttt{claude-haiku-4-5-20251001})}
\begin{quote}\small\ttfamily\raggedright
\{\\
  "audience": "Graduate students and researchers in algebraic topology, group theory, and geometric topology, particularly those interested in braid groups, configuration spaces, mapping class groups, homological algebra, or cohomological methods in discrete groups.",\\
  "key\_contributions": [\\
    "Complete characterization of K(\ensuremath{\pi},1)-spaces among configuration spaces, showing most surfaces (genus \ensuremath{\ge} 1) have this property, with explicit analysis for exceptional cases like the sphere and projective plane",\\
    "Proof that pure braid groups on surfaces are polyfree with explicit formulas for cohomological dimension in terms of genus and marked points",\\
    "Development of function space interpretation using Dold-Thom constructions relating loop spaces of configuration spaces to algebraic structures",\\
    "Explicit calculations of cohomology groups for braid and mapping class groups of low-genus surfaces, including torsion invariants",\\
    "Theory of polyfree groups establishing structural properties: finite presentation, type FL, bounded cohomological dimension, and Euler characteristic formulas",\\
    "Analysis of orbit configuration spaces under group actions and applications to symplectic groups and automorphisms of free groups"\\
  ],\\
  "plain\_language\_summary": "Configuration spaces are fundamental objects in topology that parameterize arrangements of distinct points on a surface. This work develops a comprehensive theory connecting configuration spaces to braid groups and mapping class groups---groups that describe symmetries and deformations of surfaces. The key insight is determining which configuration spaces are K(\ensuremath{\pi},1)-spaces, meaning their geometric shape is entirely determined by their fundamental group (the group of loops within the space). The authors prove that configuration spaces on most surfaces satisfy this property, extending classical results from the plane to higher-genus surfaces, with modifications needed only for the sphere and projective plane.\textbackslash{}n\textbackslash{}nThe second major theme involves function space interpretation: understanding how spaces of loops in configuration spaces relate to algebraic structures through homological methods. The third component provides explicit cohomological calculations---computing algebraic invariants that measure geometric properties---for braid and mapping class groups on surfaces of genus 0, 1, and 2, including cases with marked points. A concrete example is analyzing how k particles move on a punctured torus under equivalence by the natural SL(2,\ensuremath{\mathbb{Z}}) symmetry action.\textbackslash{}n\textbackslash{}nCentral to the work is the study of polyfree groups, which have a special hierarchical structure forming a chain of normal subgroups with free factors. The paper establishes that pure braid groups on surfaces are polyfree and determines their cohomological dimension, Euler characteristic, and related homological properties, providing tools for computing cohomology and torsion invariants in related discrete groups.",\\
  "tldr": "These notes analyze configuration spaces and braid groups by establishing their connections to K(\ensuremath{\pi},1)-spaces and computing cohomological properties of related discrete groups across surfaces of varying genus."\\
\}\\
\end{quote}

\subsection*{technical\_correctness (\texttt{claude-opus-4-7})}
\begin{quote}\small\ttfamily\raggedright
\{\\
  "claims": [\\
    \{\\
      "assessment": "supported",\\
      "claim": "Fadell-Neuwirth fibration: For a boundaryless manifold M and k \ensuremath{\ge} m \ensuremath{\ge} 1, the projection \ensuremath{\pi}\_\{k,m\}: F(M,k) \ensuremath{\to} F(M,m) is a fibration with fiber F(M \ensuremath{-} Q\_m, k\ensuremath{-}m).",\\
      "evidence": "This is a well-established classical theorem (Fadell--Neuwirth, Math. Scand. 1962). It is properly cited via reference [FN] in the bibliography and used consistently throughout the paper.",\\
      "id": "C1",\\
      "location": "Section 2, Theorem 5",\\
      "severity": "info",\\
      "suggested\_fix": null\\
    \},\\
    \{\\
      "assessment": "supported",\\
      "claim": "For M = R\^{}2 or a closed 2-manifold of genus \ensuremath{\ge} 1 (orientable or not), F(M \ensuremath{-} Q\_m, k) is a K(\ensuremath{\pi},1) for all k \ensuremath{\ge} 1, m \ensuremath{\ge} 0.",\\
      "evidence": "The induction-on-k proof using Theorem 5 (Fadell-Neuwirth) and the homotopy long exact sequence is sound. Base case rests on M \ensuremath{-} Q\_m being homotopy equivalent to either M (K(\ensuremath{\pi},1)) or a bouquet of circles (K(F,1)). The argument is standard in algebraic topology and matches the literature (Fadell--Neuwirth, see [FN]).",\\
      "id": "C2",\\
      "location": "Section 2, Theorem 7",\\
      "severity": "info",\\
      "suggested\_fix": null\\
    \},\\
    \{\\
      "assessment": "partially\_supported",\\
      "claim": "For M = R\^{}2 or a closed 2-manifold of genus g \ensuremath{\ge} 1, the pure braid group PB\_k(M \ensuremath{-} Q\_m) is polyfree with a polyfree series of length k whose i-th factor has rank 2g + m \ensuremath{-} 1 \ensuremath{-} i + k (with the convention g = 1/2 for M = R\^{}2).",\\
      "evidence": "The induction step uses the Fadell-Neuwirth fibration \ensuremath{\pi}\_\{k,k-1\} with fiber F(M \ensuremath{-} Q\_\{m+k-1\}, 1), correctly yielding a free factor of rank 2g + m \ensuremath{-} 2 + k. Combining this with the inductive series gives a polyfree series of length k. The exponent formula 2g + m \ensuremath{-} 1 \ensuremath{-} i + k matches this at the boundary values (i=k gives 2g+m\ensuremath{-}1, the rank of \ensuremath{\pi}\_1(M \ensuremath{-} Q\_m); i=1 gives 2g+m\ensuremath{-}2+k, the rank of the top fiber added in the inductive step), but the proof explicitly defers the verification of the indexing to the reader ('It is an easy exercise which will be left to the reader'). The unconventional choice g = 1/2 for R\^{}2 is a notational hack that compresses two cases but is not rigorously justified --- the underlying claim \ensuremath{\pi}\_1(R\^{}2 \ensuremath{-} Q\_m) is free of rank m needs separate treatment.",\\
      "id": "C3",\\
      "location": "Section 5 (Polyfree groups), Theorem 40",\\
      "severity": "minor",\\
      "suggested\_fix": "Provide a fully worked verification of the exponent formula in at least one nontrivial case (e.g., g=1, m=0, k=3) and explicitly handle M = R\^{}2 as a separate base case rather than via the g = 1/2 convention. Alternatively, ship a Lean/Coq formalization of the recursion at src/proofs/PolyfreeExponents.lean to mechanically verify the indexing."\\
    \},\\
    \{\\
      "assessment": "supported",\\
      "claim": "If G is polyfree with a polyfree series of length n, then (a) cd(G) \ensuremath{\le} n, (b) G is of type FL, (c) G is finitely presented, and (d) if every exponent d\_k \ensuremath{\ge} 2 then Z(G) = 1.",\\
      "evidence": "The induction on n combines (i) P. Hall's theorem on finite presentation in extensions (cited as Theorem 42, [Robin]), (ii) the LHS-spectral sequence vanishing argument for cohomological dimension (E\_2\^{}\{p,q\}=0 for p>1 or q>n\ensuremath{-}1 forces E\_2=E\_\ensuremath{\infty}), and (iii) the standard centerless argument for free groups of rank \ensuremath{\ge} 2. The reasoning is correct and reproduces classical results (Bieri, Brown, see [Brown]).",\\
      "id": "C4",\\
      "location": "Section 5, Proposition 43",\\
      "severity": "info",\\
      "suggested\_fix": null\\
    \},\\
    \{\\
      "assessment": "supported",\\
      "claim": "A finite-index subgroup H of a polyfree group G with series of length n is itself polyfree with a polyfree series of length n.",\\
      "evidence": "The proof uses Nielsen--Schreier (finite-index subgroups of finitely generated free groups are finitely generated free) and the elementary fact that intersection with a normal series gives a normal series of the subgroup. The injection H\_i/H\_\{i-1\} \ensuremath{\to} G\_i/G\_\{i-1\} into a finitely generated free group has finite index, hence H\_i/H\_\{i-1\} is free of finite rank. The argument is standard and correct.",\\
      "id": "C5",\\
      "location": "Section 5, Proposition 45",\\
      "severity": "info",\\
      "suggested\_fix": null\\
    \},\\
    \{\\
      "assessment": "partially\_supported",\\
      "claim": "For a polyfree group G with polyfree series of length n and exponents d\_k, cd(G) = n, every polyfree series of G has length n, and \ensuremath{\chi}(G) = ∏\_\{k=1\}\^{}n (1 \ensuremath{-} d\_k).",\\
      "evidence": "The upper bound cd(G) \ensuremath{\le} n comes from Proposition 43. The lower bound is obtained via Lemma 47 by constructing a finite-index subgroup T with H\^{}n(T; F\_2) \ensuremath{\ne} 0 via an LHS-spectral-sequence non-vanishing argument. The key step uses that conjugation acts trivially on H\^{}1(G\_1; F\_2) after passing to the finite-index subgroup K, so the E\_2 page splits as a tensor product. This step is correct but the spectral sequence argument relies on E\_2\^{}\{n-1,1\} surviving to E\_\ensuremath{\infty}, which requires that no nonzero differential lands in or leaves bidegree (n-1,1). The proof asserts this via the shape of E\_2, but a more explicit check of the differentials d\_2, d\_3, ... would strengthen the exposition. The Euler characteristic identity follows from Hall's identity \ensuremath{\chi}(G) = \ensuremath{\chi}(N)\ensuremath{\chi}(Q) for extensions of FP groups, which the lemma proves inductively.",\\
      "id": "C6",\\
      "location": "Section 5, Proposition 48 with Lemma 47",\\
      "severity": "minor",\\
      "suggested\_fix": "Strengthen the spectral-sequence step in Lemma 47 by explicitly noting that all higher differentials out of E\_r\^{}\{n-1,1\} or into E\_r\^{}\{n-1,1\} land in bidegrees that are zero by the shape of E\_2 (p > n-1 or q > 1). A formal verification of the LHS argument at src/proofs/PolyfreeCD.lean using a categorical spectral-sequence library would mechanically rule out subtleties."\\
    \},\\
    \{\\
      "assessment": "supported",\\
      "claim": "Stallings--Swan theorem: cd(G) \ensuremath{\le} 1 if and only if G is a free group.",\\
      "evidence": "The claim is cited as 'a deep result of Stallings and Swan' without re-proof, consistent with standard references such as Brown's 'Cohomology of Groups' (cited as [Brown]). The result is a well-established theorem in group cohomology.",\\
      "id": "C7",\\
      "location": "Section 4 (Cohomology of groups), subsection on cohomological dimension",\\
      "severity": "info",\\
      "suggested\_fix": null\\
    \},\\
    \{\\
      "assessment": "unsupported",\\
      "claim": "The mathematical content of the paper --- all theorems on configuration spaces, braid groups, polyfree groups, cohomological dimension, Euler characteristics, and LHS-spectral-sequence calculations --- is correct.",\\
      "evidence": "Per the Proof-as-Code Axiom for math.* fields: although the paper's individual results are classical or carefully argued, no executable formal verification artifact (Coq, Lean, Agda, or Isabelle/HOL) is shipped to mechanically certify the recursion arguments (Theorem 7, Theorem 40, Proposition 43, Lemma 47, Proposition 48). For a 'notes' / survey style paper that reproves classical results, the absence of a formalization repository is evidence that the load-bearing recursions have not been independently mechanically checked.",\\
      "id": "C8",\\
      "location": "Whole paper",\\
      "severity": "major",\\
      "suggested\_fix": "Ship a formalization repository with at least the core induction lemmas: (i) Fadell-Neuwirth fibration consequence for K(\ensuremath{\pi},1) (src/proofs/FNKpi1.lean), (ii) polyfree exponent recursion (src/proofs/PolyfreeExponents.lean), (iii) the LHS-spectral-sequence cohomological-dimension lower bound (src/proofs/PolyfreeCD.lean). Mathlib already contains the relevant spectral-sequence and group-cohomology infrastructure."\\
    \},\\
    \{\\
      "assessment": "partially\_supported",\\
      "claim": "The convention g = 1/2 for M = R\^{}2 in the polyfree exponent formula correctly subsumes the planar case.",\\
      "evidence": "With g = 1/2, the formula 2g + m \ensuremath{-} 1 \ensuremath{-} i + k for the exponent of the i-th factor produces m \ensuremath{-} i + k, which matches the rank of \ensuremath{\pi}\_1(R\^{}2 \ensuremath{-} Q\_\{m+k-1\}) = F\_\{m+k-2\} when one accounts for the index shift in Remark 41 (PB\_k(R\^{}2) = PB\_\{k-1\}(R\^{}2 \ensuremath{-} Q\_1)). However, the assertion 'these formulas also work for M = R\^{}2 if we set g = 1/2' is not derived; it is offered as a notational hack. The reader must verify it by hand, and the off-by-one between Theorem 40 (length k) and the planar case (length k\ensuremath{-}1, per Remark 41) is a potential source of confusion.",\\
      "id": "C9",\\
      "location": "Section 5, Theorem 40 and Remark 41",\\
      "severity": "minor",\\
      "suggested\_fix": "Rewrite Theorem 40 with an explicit case split between R\^{}2 and closed surfaces of genus g \ensuremath{\ge} 1, or alternatively prove the g = 1/2 convention as a separate lemma. State the exponent formula for PB\_k(R\^{}2) directly using Remark 41."\\
    \},\\
    \{\\
      "assessment": "unsupported",\\
      "claim": "Genus 2 mapping class group has a 'configuration-like' description from which the torsion in its cohomology follows at once, with applications to the integral cohomology of Sp(4, Z).",\\
      "evidence": "This headline claim from the Introduction is asserted but not substantiated in the portion of the paper available for review. The promise that 'torsion in the cohomology follows at once' from a configuration-like model is a strong assertion; without the explicit computation visible in the inspected text, and without an executable cohomology computation (e.g., via Sage, GAP, or a formalized Lean calculation against a CW model), the claim is at most a survey-level pointer.",\\
      "id": "C10",\\
      "location": "Introduction (and Section 15 per the TOC)",\\
      "severity": "major",\\
      "suggested\_fix": "Either provide explicit chain-level computations of H\^{}*(MCG\_2; Z) in Section 15 with explicit torsion subgroup generators, or ship a computer-algebra script (experiments/section15/mcg2\_cohomology.py using Sage's group cohomology package, or a GAP script with HAP) that reproduces the claimed torsion structure and the consequent Sp(4, Z) cohomology computation."\\
    \}\\
  ],\\
  "confidence": 0.6,\\
  "overall\_correctness": "mostly\_sound"\\
\}\\
\end{quote}

\section*{Corrections}
% corrections rendered from corrections table; empty on first publish
No corrections have been recorded.

\section*{Bibliography}
\begin{enumerate}[leftmargin=*]
  \item Bred: \{Bredon,\} Topology and Geometry, Springer Verlag GTM 139, New York-Heidelberg-Berlin, 1993.
  \item Brown: \{K. S. Brown,\} Cohomology of Groups, Springer Verlag GTM 87, New York-Heidelberg-Berlin, 1994.
  \item BM: \{Bergau, and Mennicke\}
  \item BH: \{Birman, and H. Hilden,\}
  \item C1: \{F. R. Cohen,\} On the hyperelliptic mapping class groups, \$SO(3)\$, and Spin\$\^{}c(3)\$, American J. Math., \{\textbackslash{}bf 115\} (1993), 389--434.
  \item C2: \{F.R. Cohen,\} On configuration spaces, their homology, and Lie Algebras, Journal of Pure and Applied Algebra \{\textbackslash{}bf 100\} (1995), 19-42.
  \item CG: \{F.R. Cohen, S. Gitler,\} (i) \{On loop spaces of configuration spaces\}, submitted to the Memoirs of the AMS, and (ii) \{Loop spaces of configuration spaces, braid-like groups, and knots\}, to appear in the Proceedings of the 1998 Barcelona Conference on Algebraic Topology.
  \item CLM: \{F.R. Cohen, T.J. Lada and J.P. May,\} The homology of iterated loop spaces, Lecture Notes in Math., vol. \{\textbackslash{}bf 533\}, Springer-Verlag (1976).
  \item CS: \{F.R. Cohen, T. Sato,\} in preparation.
  \item CX: \{F.R. Cohen, M. Xicot\textbackslash{}'encatl,\} submitted.
  \item DC: \{D. C. Cohen,\} Monodromy of fibre-type arrangements and orbit configuration spaces, to appear in Forum Mathematicum.
  \item Dan: \{D. C. Cohen, Alexander I. Suciu,\} Homology of iterated semidirect products of free groups, J. Pure Appl. Algebra \$\textbackslash{}bf\{126\}\$ (1998), pg 87-120.
  \item EE: \{C. Earle, and J. Eells,\} A fibre bundle description of Teichm\textbackslash{}"uller theory, J. Differential Geometry \{\textbackslash{}bf 3\}(1969), 19-43.
  \item FH1: \{E. Fadell and S. Husseini,\} The space of loops on configuration spaces and the Majer-Terracini index, Topol. Methods in Nonlinear Anal., Journal of the Julius Schauder Center, \{\textbackslash{}bf 11\} (1998), 249--271.
  \item FH2: \{E. Fadell and S. Husseini,\} Geometry and Topology of Configuration Spaces, in preparation.
  \item FN: \{E. Fadell and L. Neuwirth,\} Configuration spaces, Math. Scand. \{\textbackslash{}bf 10\} (1962), 119-126.
  \item FR: \{M. Falk, and R. Randell,\} The lower central series of a fiber-type arrangement, Invent. Math. \{\textbackslash{}bf 82\} (1985), 77-88.
  \item GP: \{V. Guillemin, and A. Pollack,\} Differential Topology, Prentice-Hall, 1974.
  \item H: \{R. Hain,\} Infinitesimal presentations of the Torelli group, J. Amer. Math. Soc. \{\textbackslash{}bf 10\} (1997),597-691.
  \item K1: \{T. Kohno,\} Linear represenations of braid groups and classical Yang-Baxter equations, Cont. Math. \{\textbackslash{}bf 78\} (1988), 339-363.
  \item K2: \{T. Kohno,\} Vassiliev invariants and de Rham complex on the space of knots, Cont. Math. \{\textbackslash{}bf 179\} (1994), 123-138.
  \item K3: \{T. Kohno,\} Elliptic KZ system, braid groups of the torus and Vassiliev invariants, Topology and its Applications, \{\textbackslash{}bf 78\} (1997), 79-94.
  \item KO: \{T. Kohno, and T. Oda,\} The lower central series of the pure braid group of an algebraic curve, Advanced Studies in Pure Mathematics, Galois Represntations and Arithmetic Algebraic Geometry \{\textbackslash{}bf 12\} (1987), 201-219.
  \item MM: \{J. W. Milnor, and J.C. Moore,\} On the structure of Hopf algebras, Ann. of Math. \{\textbackslash{}bf 81\} (1965), 211-264.
  \item Robin: \{Derek J.S. Robinson,\} A Course in the Theory of Groups, Springer Verlag GTM 80, New York-Heidelberg-Berlin, 1995.
  \item S: \{N. E. Steenrod,\} The Topology of Fibre Bundles, Princeton University Press.
  \item Sa: \{T. Sato,\} On free groups and morphisms of coalgebras, student topology seminar, Univ. of Roch., Spring 1998, preprint.
  \item Smale: \{S. Smale,\} P.A.M.S.
  \item Span: \{Edwin H. Spanier,\} Algebraic Topology, Springer Verlag, New York-Heidelberg-Berlin, 1966.
  \item White: \{Whitehead,\} Elements of Homotopy Theory, Springer Verlag GTM 61, New York-Heidelberg-Berlin, .
  \item W: \{S. Willerton,\} thesis.
  \item X: \{M. Xicot\textbackslash{}'encatl,\} Orbit configuration spaces, infinitesimal braid relations, and equivariant function spaces, Ph.D. thesis, University of Rochester, Spring 1997.
  \item Xia: \{Yining Xia,\}
\end{enumerate}

% \bibliographystyle{plain}
% \bibliography{references}
\end{document}
